\documentclass[10pt,a4paper]{amsart}
\usepackage[margin=19mm]{geometry}
\usepackage{amsmath,amssymb,mathtools}
\usepackage[scr=rsfs,cal=euler]{mathalfa}
\usepackage{xfrac}
\usepackage[unicode,hidelinks,bookmarks=false]{hyperref}
\newcommand{\ie}{i.e.\ }
\newcommand{\cf}{cf.\ }
\newcommand{\resp}{respectively\ }
\newcommand{\Subsection}{\subsection}
\providecommand{\nobreakdash}{\nobreak\hbox{-}}
\setlength{\emergencystretch}{2em}
\multlinegap0pt
\usepackage{bm}
\usepackage{bbm}
\usepackage{dsfont}
\usepackage{xspace}
\usepackage{cancel}
\usepackage{mathtools}
\usepackage{amsthm}
\makeatletter
\@ifundefined{theorem}{\newtheorem{theorem}{Theorem}}{}
\makeatother
\title[Tight Bounds for Cycle-Edge Decompositions and Covers]{Tight Bounds for Cycle-Edge Decompositions and Covers}
\author{Saieed Akbari, Jonny Aloni, Arash Beikmohammadi, Alexander Clow}
\date{Source: arXiv:2509.01901v2 (CC BY 4.0)}
\begin{document}
\maketitle

\noindent\emph{Source and licence.} Tight Bounds for Cycle-Edge Decompositions and Covers. Version arXiv:2509.01901v2 (CC BY 4.0), distributed under the Creative Commons Attribution 4.0 International licence, as recorded in the arXiv licence metadata for this exact version and shown on the abstract page https://arxiv.org/abs/2509.01901v2. Full rights notice: licenses/arxiv-2509.01901-CC-BY-4.0.txt.\par\smallskip

\noindent\textit{Editorial scope.} Counted unit: the Theorem whose source label is \texttt{Thm: n-3 even} in the article (source0.tex lines 357--363, source label \texttt{Thm: n-3 even}), delivered together with the article's own complete original proof of it (source0.tex lines 364--447, one \texttt{proof} environment of 84 lines). No mathematics is written, repaired or re-derived: statement and proof are the article's own text line for line. Required context retained uncounted: Thm: Eulerian decomp (lines 180--182). Every definition, display and label of those lines is retained unchanged. Omitted from this package and not redistributed: 1--179, 183--356, 448--844. Nothing inside a retained excerpt is omitted or altered: every retained body is delivered exactly as the article's lines are written, so no omission receipt of any kind accompanies this package. Citations: the retained excerpts carry no citation command at all, so no bibliography entry of the article is transcribed and no reference is redistributed. Reference adaptation: none was needed and none was made; every reference inside the delivered unit resolves inside it, so no unresolved reference or citation is produced. Printed slips kept literal: no printed word, symbol, hypothesis or number of the counted statement or of its proof was altered. Layout and class: the article's own class is \texttt{amsart} with packages including hyperref, amsfonts, graphicx, tabularx, array, color, comment, amsmath, amsthm, amssymb, fullpage, xcolor, tikz, listings; the wrapper is the lane's pinned \texttt{amsart} editorial wrapper at 10pt on A4, which restates the theorem-like environments the retained text needs and adds the article's own macro definitions that the retained text needs (none), with extra wrapper packages (bm, bbm, dsfont, xspace, cancel, mathtools). The delivered numbering is the unit's own. Assets: no retained span contains a figure environment, a graphics-inclusion command, a PDF-page inclusion, a table or a TikZ picture; no image file is copied into this package, the package declares no asset dependency and no image file is redistributed with it. Licence: version arXiv:2509.01901v2 of this article is distributed under the Creative Commons Attribution 4.0 International licence, as recorded in the arXiv licence metadata for this exact version and shown on the abstract page https://arxiv.org/abs/2509.01901v2. A full rights notice is delivered at licenses/arxiv-2509.01901-CC-BY-4.0.txt. Characters: every retained excerpt is pure ASCII, so the delivered engine, XeLaTeX with its default Latin Modern text font, reports no missing character.
\begin{theorem}\label{Thm: Eulerian decomp}
    Every $n$ vertex graph containing a cycle can be decomposed into an even graph and at most $n-2$ edges.
\end{theorem}

\begin{theorem}\label{Thm: n-3 even}
Let $G = (V,E)$ be a simple connected graph that contains a cycle. Then exactly one of the following are true
\begin{itemize}
\item[i)] $G$ is decomposable by an even subgraph and at most $n-3$ edges.
\item[ii)] $G$ is $K_4$ with a disjoint tree hanging off every vertex of the $K_4$.
\end{itemize}
\end{theorem}

\begin{proof}
By Theorem~\ref{Thm: Eulerian decomp} $G$ can be decomposed into an even graph and $n-2$ edges. 
If $G$ is a graph of type ii), then any even graph in such a decomposition will be of size at most 4 and thus an additional $n-2$ edges are needed in the decomposition.
So it remains to show that all graphs that contain a cycle which cannot be decomposed by $n-3$ edges are of type ii).

Let $G$ be an $n$ vertex graph that contains a cycle.
Let $H$ be an edge maximum even graph of $G$.
Then $F:=G - E(H)$ is a forest. 
By Theorem~\ref{Thm: Eulerian decomp} $F$ is not a tree, 
since this would imply $F$ has $n-1$ edges, 
thereby implying that $G$ cannot be decomposed into an even subgraph and $n-2$ edges, a contradiction.
If $F$ has at least $3$ connected components then it can be decomposed by $n-3$ edges. 
Suppose then that $F$ is the union of 2 disjoint trees $T_1$ and $T_2$. 

Note that for every edge $uv\in E(H)$ vertices $u$ and $v$ cannot be in the same component of $F$ denoted as $T_i$.
To see this, notice that if $u$ and $v$ are in $T_i$, then
there is a path $P$ in $ T_i$ that joins $u$ to $v$.
In this case $(H\cup P) - \{uv\}$ is a larger even graph than $H$.


So $H$ is a subgraph of the bipartite graph $(V,E(T_1,T_2))$,
where $E(T_1,T_2)$ is the set of edges with one endpoint in $T_1$ and the other in $T_2$.
Suppose that $H$ contains a matching of size $2$,
$u_1v_1,~u_2v_2\in E(H)$,
where $v_1,v_2\in V(T_1)$ and $u_1,u_2\in V(T_2)$.
Let $P_1$ be the path in $T_1$ that connects $v_1$ to $v_2$ and $P_2$ be the path that connects $u_1$ to $u_2$. 
The following graph is even $H^\prime:= (H\cup P_1\cup P_2) - \{u_1v_1,~u_2v_2\}$.
Then the maximality of $H$ implies $|E(H^\prime)| \leq |E(H)|$.
Hence, $v_1$ and $v_2$ are adjacent to each other in $T_1$, and $u_1$ and $u_2$ are adjacent in $T_2$. 

Suppose $H$ contains a matching of size $3$, call it $u_1v_1,u_2v_2,u_3v_3$.
Since $u_1,u_2,v_1,v_2$ in the last paragraph were chosen without loss of generality, we conclude that the maximality of $H$ implies $v_1,v_2,v_3$ and $u_1,u_2,u_3$ induce triangles in $F$.
But this is a contradiction since $F$ is a forest.
Hence, $H$ has no matching of size $3$.


Since $H$ is bipartite, K\H{o}nig's theorem implies the matching number of $H$, which is at most $2$, 
is equal to the vertex cover number of $H$.
Hence, the vertex cover number of $H$ is at most $2$.
Notice the vertex cover number
of $H$ must be greater than $1$, since $H$ contains a cycle by our assumption that $G$ contains a cycle.
Moreover, no connected component of $H$ is a star since $H$ is 2-regular. 
Therefore, $H$ has vertex cover number $2$, and $H$ has exactly one connected component containing $2$ or more vertices.
Suppose this component is $Q$.


Let $\{a,b\}\subseteq V$ be a vertex cover in $Q$.
We consider the cases 
$a \in V(T_1)$ and $b\in V(T_2)$, versus $a,b\in V(T_1)$ separately.
The labels for vertices and $T_1$ and $T_2$ are chosen without loss of generality.


\medskip
\noindent \underline{Case 1:} $a \in V(T_1)$ and $b\in V(T_2)$.
\medskip

Since the matching number of $Q$ is $2$, $\{b\}$ is not a vertex cover.
Thus there exists a vertex $u\in V(T_2)$ which is not $b$ and $ua\in E(H)$. The degree of $u$ is 1 because $u$ can only be adjacent to elements in the vertex cover $\{a,b\}$ and $u$ is not adjacent to $b$ because they are both in $T_2$.
This contradicts $H$ being even. So this case does not occur.

\medskip
\noindent \underline{Case 2:} $a,b \in V(T_1)$.
\medskip

Let $k = |V(T_2)\cap V(Q)|$.
Since $\{a,b\}$ is a vertex cover, all edges are incident to $a$ or $b$.
As $a,b \in V(T_1)$, the fact that $Q$ is bipartite implies $V(T_1)\cap V(Q) = \{a,b\}$.
Since $Q$ has matching number $2$
both vertices $a$ and $b$ have strictly positive degree.
Since $Q$ is even and bipartite, this implies $k\geq 2$.
If $k>2$ then there exists vertices $u_1,u_2,u_3$ in $T_2$.
Each of these vertices has a positive and even degree in $Q$.
Since $Q$ is bipartite, this implies $a$ and $b$ are both neighbours of each vertex $u_i$.

As with earlier in the proof, this implies $u_1,u_2,u_3$ forms a triangle in $G$.
Since $u_1,u_2,u_3$ are in $V(T_2)$, every edge of this triangle is in $T_2$, contradicting that $T_2$ is a tree.
So $k=2$.

Since $k=2$, $V(T_2)\cap Q = \{u_1,u_2\}$.
This forces $Q$ to have $4$ vertices.
Since $V(Q)$ induces a clique, $V(Q)$ induces $K_4$.
All other edges in $G$ are edges in $T$, a tree.
From here it is trivial to verify $G$ is of type ii).
\end{proof}

\end{document}
